\documentclass[10pt,a4paper]{amsart}
\usepackage[margin=19mm]{geometry}
\usepackage{amsmath,amssymb,mathtools}
\usepackage[scr=rsfs,cal=euler]{mathalfa}
\usepackage{xfrac}
\usepackage[unicode,hidelinks,bookmarks=false]{hyperref}
\newcommand{\ie}{i.e.\ }
\newcommand{\cf}{cf.\ }
\newcommand{\resp}{respectively\ }
\newcommand{\Subsection}{\subsection}
\providecommand{\nobreakdash}{\nobreak\hbox{-}}
\setlength{\emergencystretch}{2em}
\multlinegap0pt
\usepackage{amssymb}
\usepackage[shortlabels]{enumitem}
\usepackage{tikz}
\usepackage[normalem]{ulem}
\usepackage{bm}
\usepackage{xcolor}
\usepackage{mathtools}
\newtheorem{prop}{Proposition}
\newcommand{\cD}{\mathcal{D}}
\newcommand{\cR}{\mathcal{R}}
\newcommand{\coC}{\mathbf{C}}
\newcommand{\dx}{\partial_x}
\newcommand{\dxxx}{\partial_{xxx}}
\title{From Spectral Curves to Variational Equations: Geometry and Differential Galois Theory of KdV Cnoidal Waves}
\author{Juan Jose Morales-Ruiz, Jean Pierre Ramis, Maria-Angeles Zurro}
\date{Source: arXiv:2608.25227v1 (CC BY 4.0)}
\begin{document}
\maketitle

\noindent\emph{Source and licence.} From Spectral Curves to Variational Equations: Geometry and Differential Galois Theory of KdV Cnoidal Waves. Version arXiv:2608.25227v1 (CC BY 4.0), distributed by arXiv under the Creative Commons Attribution 4.0 International licence, http://creativecommons.org/licenses/by/4.0/, as recorded in the arXiv licence metadata for this exact version and shown on the abstract page https://arxiv.org/abs/2608.25227v1. Full licence notice: \texttt{licenses/arxiv-2608.25227-CC-BY-4.0.txt}.\par\smallskip

\noindent\textit{Editorial scope.} Counted unit: the article's prop def-Q-ell (source0.tex lines 1256--1262). The article's complete original proof of it is delivered with it (source0.tex lines 1264--1298, one \texttt{proof} environment of 35 lines). No mathematics is written, repaired or re-derived: the statement and the proof are the article's own lines, in the article's own notation, and no retained line is substituted or re-flowed.  Retained uncounted as the source-local context the proof needs: lines 1038--1040.  Nothing was omitted from any retained span, so the package carries no omission record.  No retained line contains a citation, so the wrapper carries no bibliography and no bibliography entry.  No retained line contains a figure environment, an image-inclusion command or drawing code, so no image is delivered and the package declares no asset dependency.  Editorial formatting only, in addition: an AMS \texttt{amsart} preamble that restates the article's own declarations and macros used by the retained lines, namely the environments \texttt{prop} on separate counters, together with the standard packages those lines need.  The retained environments print in this document as Proposition 1 in order of appearance, whereas the article prints the same items with the numbering of its own full document.  The complete native source of the article as distributed in the arXiv e-print for this exact version is bundled unchanged as \texttt{sources/208558/source0.tex}.
 \begin{equation}\label{eq-master-densities}
-\psi_{zzz}+4 (u-E)\psi_z +2u_z \psi  =0 \quad .
\end{equation}

\begin{prop}
    Let $f$ be  a constant coefficients polynomial $f(E)=\sum_{k=0}^\ell d_k E^k $ in $\coC[E]$. Then
    \begin{equation}\label{def-Q-ell}
        Q_{\ell , f } (E) = f_{[0]}(E) v_0 +\cdots + f_{[\ell]}(E) v_\ell
    \end{equation}
    is a solution of \eqref{eq-master-densities} in $\coC\{u\}[E]$. Furthermore, any polynomial solution of \eqref{eq-master-densities} is of the form $Q_{\ell , f } $ for some $f  $ in $\coC[E]$.
\end{prop}

\begin{proof}
    Let $Q_\ell =\sum_{k=0} ^\ell \omega_k E^k$ be a polynomial solution of \eqref{eq-master-densities} in $\coC\{u\}[E]$ of degree $\ell >0$. Then equating degrees in $E$ we obtain the differential recursion
    \begin{equation}\label{recursion-D2}
        -\cD_2 (\omega_0 )=0 \ , \ -\cD_2 (\omega_n )=4 \partial (\omega_{n-1}) \ , \ 4 \partial (\omega_\ell )=0 \ .
    \end{equation}
    with $ \cD_2 = \dxxx - 4u\,\dx - 2u_x $. 

For $ Q_{\ell , f } (E) = f_{[0]}(E) v_0 +\cdots + f_{[\ell]}(E) v_\ell$ to be a solution we rewrite it as
\[
 Q_{\ell , f } (E) = \sum_{n=0}^{\ell} \left( \sum_{k=0}^{\ell -n} d_{n+k}  v_k \right) E^n  = \sum_{n=0}^{\ell}  \omega_n  E^n 
\]
with $\omega_n = \sum_{k=0}^{\ell -n} d_{n+k}  v_k$. Then the sequence $\omega_0 , \dots , \omega_\ell$ satisfies the differential recursion \eqref{recursion-D2}, which implies that $ Q_{\ell , f } $ is a solution of \eqref{eq-master-densities}.

Next we proceed by induction on $n$ to prove that we can construct a finite sequence $c_\ell , \dots c_n$ of constants such that $Q_\ell =\sum_{k=0} ^\ell \omega_k E^k=   Q_{\ell , f }$ for $f(E)=\sum_{n=0}^\ell c_n E^n$. In fact, since $\partial (\omega_\ell )=0$, we have $\omega_\ell =c_\ell$ for some constant $c_\ell \in \coC$. Consequently $4 \partial (\omega_{\ell -1}) = -\cD_2 (c _\ell  )=2u_z c_\ell$. Thus, $\omega_{\ell -1}= c_\ell v_1 +c_{\ell -1} v_0 $ for some constant $c_{\ell -1} \in \coC $. Suppose now that the formula
    \[
    \omega_{n+1} = \sum_{j=0}^{\ell-n-1}c_{n+1+j} v_j \ 
    \]
    holds, and let us prove it for $n$ by induction. We use recursion \eqref{recursion-D2} for this purpose.
    \[
    4 \partial (\omega_{n})= -\cD_2 (\omega_{n+1} )=4\partial \cR^* (\omega_{n+1})=4\partial 
    \left(
    \sum_{j=0}^{\ell-n-1} c_{n+1+j} v_{j+1}
    \right) .
    \]
 Therefore, $\omega_{n} = \sum_{k=1}^{\ell-n} c_{n+k} v_k + c_n v_0 $ for some constant $c_n \in \coC$. Consequently, we can write $Q_\ell $ as $Q_{\ell , f}$ with $f(E) = \sum_{k=0}^\ell c_k E^k$ because
   \[
   Q_{\ell } (E) = \sum_{n=0}^{\ell} \sum_{k=0}^{\ell -n} c_{n+k}  E^n  v_k= \sum_{k=0}^{\ell} 
   \left(
   \sum_{j=0}^{\ell -k} c_{k+j}E^j 
   \right)
   v_k = \sum_{k=0}^{\ell} f_{[k]} (E) 
     v_k \ ,
   \]
  and we achieve the statement.
\end{proof}

\end{document}
